\chapter{Initial Training Programs in Exa-MA Domains}

\section{Overview of Initial Training Programs}
Brief description of training programs in Exa-MA domains of math and computer science.
\subsection{Data Collection and Reporting}
Explanation of data collection methodology and community engagement.
\subsubsection{Form Design and Distribution}
Summary of the form provided to community members for reporting relevant training programs.

possible form format:
\begin{itemize}
    \item \textbf{Program Name}: (e.g., Master in Applied Mathematics)
    \item \textbf{Institution}: (e.g., University, Engineering School)
    \item \textbf{WP Relevance}: (checkboxes or multi-select for WP1 to WP6)
    \item \textbf{Relevant Courses/Modules}: (e.g., Numerical Methods, Data Science)
    \item \textbf{Target Level}: (e.g., Master’s, Engineering, PhD)
    \item \textbf{Primary Focus}: (e.g., theory, application, software training)
    \item \textbf{Contact Person/Instructor} (optional)
\end{itemize}

should we require that training programs support hybrid computing and advanced methods or is it sufficien to have an introductory course in HPC and numerical methods?

\subsubsection{Compilation of Initial Training Programs}
Table or list of initial training programs as reported by community members.
\subsection{Analysis of Training Program Alignment}
Analysis of how identified programs support Exa-MA's goals.

\section{Summary of Program Contributions by WP}
\subsection{WP1 - Discretization}
\subsection{WP2 - Model Order, Surrogate, Scientific Machine Learning Methods}
\subsection{WP3 - Solvers}
\subsection{WP4 - Data Assimilation}
\subsection{WP5 - Optimization}
\subsection{WP6 - Uncertainty Quantification}

\